\chapter{Introduction}

\section{Background and Context}

Cloud storage separates durable file retention from the physical limits of a personal device, but conventional services retain files without helping users discover information when the exact filename or wording is unknown. NimbusVault is a storage and document-intelligence platform that combines secure file organisation, controlled collaboration and document intelligence. A single permission policy governs downloads, search results, retrieved evidence and saved conversations.

The platform separates object storage from relational ownership, and asynchronous processing from heavy inference. BM25 ranks lexical matches \cite{RZ09}; sentence embeddings capture meaning \cite{RG19,WWD20,STMiniLM}; and retrieval-augmented generation (RAG) conditions answers on retrieved evidence \cite{LPP20}. The contribution of this project is the integration of these techniques with current permissions, evidence validation and reproducible testing, rather than a new language model or retrieval algorithm.

NimbusVault is built as a web application. A React interface talks to an Express server that owns identity, folders, ownership and sharing through PostgreSQL. Original files are kept in private Amazon S3. A private FastAPI service extracts text, builds embeddings and answers questions with a locally hosted Qwen model, using only the files the user is allowed to read.

\section{Motivation}

In the Software Engineering course we had to choose a real problem and carry it from proposal through design, incremental implementation and testing. Cloud storage is something every student uses, and finding a document later usually depends on remembering its name or exact wording. We wanted a project where search and question answering work over stored files without weakening the rules about who can see what.

The project also let us work with real services instead of a hard-coded demo: private S3 storage, an event-driven AWS workflow, and a language model hosted on our own hardware rather than a paid API. The original proposal was refined during development to Express, PostgreSQL, SQLite and a self-hosted Qwen model. No commercial-product comparison or distributed availability benchmark was performed.

\section{Problem Statement}

Cloud storage by itself does not solve document discovery, controlled sharing or evidence-grounded question answering. A user who cannot remember a filename or exact wording cannot find the document. An answer generated over stored documents is only trustworthy if it is limited to what that user may read and is backed by checkable evidence. The project therefore addresses how to provide secure file storage together with semantic search and question answering, in which every download, search result, citation and saved conversation obeys one permission policy, and in which insufficient evidence leads to abstention rather than an unsupported answer.

\section{Project Objectives}
\label{sec:objectives}

The status of each objective is stated as it stands in the implementation, not as it was planned.

\begin{enumerate}
  \item \textbf{Secure cloud storage:} Keep original files in private Amazon S3, with historical synthetic upload and SHA-256 readback evidence.
  \item \textbf{File management:} Provide nested folders, upload, rename, move, trash and restore through PostgreSQL-backed APIs.
  \item \textbf{Controlled sharing:} Support resource grants, inherited access and bounded share links.
  \item \textbf{Intelligent discovery:} Offer MiniLM semantic retrieval, SQLite FTS5 BM25 and hybrid ranking over the files a user may read.
  \item \textbf{Document question answering:} Generate answers with a local Qwen model using validated citations and explicit abstention.
  \item \textbf{Deployment readiness:} Run as a standalone Windows application and in a native Linux rehearsal; complete EC2 deployment is pending.
\end{enumerate}

Each objective is paired with the design in Chapter~\ref{ch:design}, the implementation in Chapter~\ref{ch:impl} and the recorded evidence in Chapter~\ref{ch:results}.
